% =============================================================================
% Patch: §4 Finite-scale local blocks and whitening (hardened)
%   sec:finite-scale-local-blocks-and-whitening
%
% Target file:   rh/rh_rebuild.tex
% Target lines:  919 through 1004 (current §4 stub)
%
% Action: this patch is NOT to be merged into rh_rebuild.tex yet.  It is
% staged under rh/patches/ for review.
%
% Migration content (relative to the current stub), with hardening from
% the referee diff at
%   rh/patches/referee/sec-finite-scale-local-blocks-and-whitening.hardened.tex:
%
%   - Definition def:finite-packet-domain: introduce the finite packet
%     domain D_T = {(m, s) : t_± in I_T} with signed s, so the s -> 0
%     removable-singularity statement is precise.
%
%   - Lemma lem:finite-blocks-exact-interpretation (promoted from a
%     remark): G_{m,±}(s) = J(t_±) and N_m(s) is the kernel-derivative
%     form at (T_1, T_2) = (t_-, t_+).  Exact, citable downstream.
%
%   - Lemma lem:Nm-removable-singularity: continuous removable
%     singularity N_m(0) = J(m).  Proof rewritten with correct
%     o-expansions:
%         Delta_m(s) = -q s - q'' s^3 / 24 + o(s^3),
%         q_+ = q + (s/2) q' + (s^2/8) q'' + o(s^2),
%         q_- = q - (s/2) q' + (s^2/8) q'' + o(s^2).
%     (Earlier version had the q_± linear-term sign flipped in the
%     proof prose; final limits were right but the proof line was wrong.)
%
%   - Remark rem:Nm-coalescence-not-cross-limit: the s -> 0 coalescence
%     is NOT obtained from lem:cross-block-jet-limit (which is fixed-
%     separation in h, not in s).  It is the separate Taylor
%     computation above.
%
%   - Definition def:whitened-finite-block: explicit "positive definite
%     inverse square root" reference.
%
%   - Corollary cor:omega-coalescence: Omega_zeta_hat(0; m) = I_2.
%     Proof now cites continuity of the SPD square-root map.
%
%   - Lemma lem:whitening-loss: operator-norm bound with explicit
%     epsilon_T -> 0 uniformly on D_T; C_W = 0 derived from
%     lambda_min(J(t)) -> infty uniformly for t in I_T after raising
%     the cutoff.
%
%   - Remark rem:whitening-loss-scope: scope sharpened.  The lemma
%     controls only real-axis same-point whitening for the theta phase;
%     no actual-zeta suppression, no source-energy, no complex-disk
%     whitening, no holomorphic estimates.
%
% Closes:
%   rem:wip-whitening-loss
%
% Status (when merged):
%   \statusmig      not-started -> migrated
%   \statusreferee  not-started -> in-progress
%   \statussympy    not-started -> verified
%       [rh/sympy/sec-finite-scale-local-blocks-and-whitening.py]
%   \statussim      not-started -> verified
%       [rh/simulation/sec-finite-scale-local-blocks-and-whitening.py]
%   \statuslean     not-started -> n/a
%
% Verification:
%   - rh/sympy/sec-finite-scale-local-blocks-and-whitening.py
%       verify_G_matches_J,
%       verify_Nm_matches_cross_block,
%       verify_q_pm_taylor_signs,
%       verify_Nm_removable_singularity,
%       verify_omega_coalescence,
%       verify_whitening_loss_constant,
%       verify_inverse_sqrt_closed_form.
%   - rh/simulation/sec-finite-scale-local-blocks-and-whitening.py
%       check_finite_blocks_at_T,
%       check_q_pm_signs_numerical,
%       check_Nm_removable_singularity_numerical,
%       check_signed_s_domain,
%       check_omega_coalescence_numerical,
%       check_whitening_loss_sweep,
%       check_lambda_min_polynomial_floor,
%       check_inverse_sqrt_routes.
% =============================================================================

% =============================================================================
\section{Finite-scale local blocks and whitening}
\label{sec:finite-scale-local-blocks-and-whitening}
\statuspaper{LIVE}\quad
\statusmig{migrated}\quad
\statusreferee{in-progress}\quad
\statussympy[rh/sympy/sec-finite-scale-local-blocks-and-whitening.py]{verified}\quad
\statussim[rh/simulation/sec-finite-scale-local-blocks-and-whitening.py]{verified}\quad
\statuslean{n/a}
\par\medskip

Around a midpoint \(m\) at finite separation \(s\), the same-point and
cross-block jets of Section~\ref{sec:jet-limit-local-blocks} give two
same-point Gram blocks \(G_{m,\pm}(s)\) and a mixed block \(N_m(s)\).
This section records the finite real-axis block package and the
whitening-factor loss used by later comparison arguments.  No
source-energy transfer or actual-zeta transverse suppression input is
used here.

\begin{definition}[Finite packet domain]
\label{def:finite-packet-domain}
For a retained theta chart \(I_T\), set \(t_\pm = m\pm s/2\) and
\[
        \mathcal D_T
        :=\bigl\{(m,s)\in\mathbb R^2 : t_-\in I_T,\ t_+\in I_T\bigr\}.
\]
Write \(\mathcal D_T^\times := \{(m,s)\in\mathcal D_T : s\ne 0\}\).
The orientation convention is that rows correspond to the center
\(t_-\) and columns to the center \(t_+\).  In applications \(s>0\),
but allowing signed \(s\) makes the removable-singularity statement at
\(s=0\) precise.  Uniform estimates below are taken on
\(\mathcal D_T\), after the global lower-height cutoff has been
enlarged if necessary.
\end{definition}

\subsection{Finite same-point and mixed blocks}
\label{subsec:finite-same-point-and-mixed-blocks}

For \((m,s)\in\mathcal D_T\), write
\[
        t_\pm = m\pm s/2,
        \qquad q_\pm = q(t_\pm),
        \qquad \Del_m(s) = \Ph(t_-)-\Ph(t_+).
\]

\begin{definition}[Finite same-point and mixed blocks]
\label{def:finite-same-point-and-mixed-blocks}
For \((m,s)\in\mathcal D_T\), the finite same-point blocks are
\begin{equation}
\label{eq:Gpm}
G_{m,\pm}(s)
=\frac1\pi
\begin{pmatrix}
2q(t_\pm) & \tfrac12 q'(t_\pm)\\[0.8ex]
\tfrac12 q'(t_\pm) & \tfrac1{12}\bigl(q''(t_\pm)+2q(t_\pm)^3\bigr)
\end{pmatrix}.
\end{equation}
For \(s\ne 0\), the finite mixed block is
\begin{equation}
\label{eq:Nm}
N_m(s)
=\frac1\pi
\begin{pmatrix}
-\dfrac{2\sin\Del_m(s)}{s}
&
\dfrac{\sin\Del_m(s)+q_+ s\cos\Del_m(s)}{s^2}
\\[2.4ex]
-\dfrac{\sin\Del_m(s)+q_- s\cos\Del_m(s)}{s^2}
&
\dfrac{(q_-+q_+)s\cos\Del_m(s)+(2-q_- q_+ s^2)\sin\Del_m(s)}{2s^3}
\end{pmatrix}.
\end{equation}
\end{definition}

\begin{lemma}[Exact finite-block interpretation]
\label{lem:finite-blocks-exact-interpretation}
For \((m,s)\in\mathcal D_T\),
\[
        G_{m,\pm}(s) = J(t_\pm).
\]
For \((m,s)\in\mathcal D_T^\times\), moreover,
\begin{equation}
\label{eq:Nm-kernel-derivative-form}
N_m(s)
=
\begin{pmatrix}
2K_\Ph(t_-,t_+) & \partial_y K_\Ph(t_-,t_+)\\[0.6ex]
\partial_x K_\Ph(t_-,t_+) & \tfrac12\partial_{xy}K_\Ph(t_-,t_+)
\end{pmatrix}.
\end{equation}
Equivalently, \(N_m(s) = \pi^{-1} N_{12}\) for the matrix \(N_{12}\) of
\eqref{eq:cross-N12} taken at \((T_1, T_2) = (t_-, t_+)\), with
\(s_{12} = T_1 - T_2 = -s\) absorbed into the displayed signs in
\eqref{eq:Nm}.
\end{lemma}

\begin{proof}
The identity \(G_{m,\pm}(s) = J(t_\pm)\) is immediate by comparing
\eqref{eq:Gpm} with \eqref{eq:same-point-J}.  For the mixed block,
apply the off-diagonal kernel-derivative formulas of
Lemma~\ref{lem:phase-kernel-derivatives} at
\((T_1, T_2) = (t_-, t_+)\); there \(s_{12} = t_- - t_+ = -s\),
\(q_1 = q_-\), \(q_2 = q_+\), and \(\Delta = \Del_m(s)\).
Substitution gives the four entries of
\eqref{eq:Nm-kernel-derivative-form}, which match \eqref{eq:Nm} entry
by entry.
\end{proof}

\begin{lemma}[Removable singularity of \(N_m(s)\)]
\label{lem:Nm-removable-singularity}
Let \(m\in I_T\).  If \(\Ph\in C^3\) on a neighborhood of \(m\), then,
as \(s\to0\) with \((m,s)\in\mathcal D_T^\times\), the entries of
\(N_m(s)\) of \eqref{eq:Nm} extend continuously to \(s=0\) with
\begin{equation}
\label{eq:Nm-at-zero}
        N_m(0) = J(m)
        = \frac1\pi
        \begin{pmatrix}
        2q(m) & \tfrac12 q'(m)\\[0.8ex]
        \tfrac12 q'(m) & \tfrac1{12}\bigl(q''(m)+2q(m)^3\bigr)
        \end{pmatrix}.
\end{equation}
\end{lemma}

\begin{proof}
Write \(q = q(m)\), \(a = q'(m)\), \(b = q''(m)\).  Taylor expansion
about \(m\), using the symmetric points \(m\pm s/2\), gives
\[
        \Del_m(s)
        = \Ph(m-s/2)-\Ph(m+s/2)
        = -qs - \frac{b}{24}s^3 + o(s^3),
\]
\[
        q_+ = q + \frac{a}{2}s + \frac{b}{8}s^2 + o(s^2),
        \qquad
        q_- = q - \frac{a}{2}s + \frac{b}{8}s^2 + o(s^2).
\]
Hence
\[
        \sin\Del_m(s)
        = -qs + \!\left(\frac{q^3}{6}-\frac{b}{24}\right)\!s^3
          + o(s^3),
        \qquad
        \cos\Del_m(s) = 1 - \frac{q^2}{2}s^2 + o(s^2).
\]
Substituting into \eqref{eq:Nm} entry by entry,
\begin{align*}
        -\frac{2\sin\Del_m(s)}{s}
        &= 2q + o(1),\\
        \frac{\sin\Del_m(s)+q_+ s\cos\Del_m(s)}{s^2}
        &= \frac{a}{2} + o(1),\\
        -\frac{\sin\Del_m(s)+q_- s\cos\Del_m(s)}{s^2}
        &= \frac{a}{2} + o(1),\\
        \frac{(q_-+q_+)s\cos\Del_m(s)+(2-q_- q_+ s^2)\sin\Del_m(s)}{2s^3}
        &= \frac{b}{12}+\frac{q^3}{6}+o(1).
\end{align*}
The common factor \(1/\pi\) in \eqref{eq:Nm} gives
\eqref{eq:Nm-at-zero}.
\end{proof}

\begin{remark}[Coalescence is not the cross-block jet limit]
\label{rem:Nm-coalescence-not-cross-limit}
Lemma~\ref{lem:cross-block-jet-limit} is a fixed-center-separation
limit in the sampling parameter \(h\); its constants are not uniform
as \(s\to 0\) (Remark~\ref{rem:cross-block-fixed-separation-scope}).
Lemma~\ref{lem:Nm-removable-singularity} is a separate coalescence
statement in the center-separation parameter \(s\), proved by direct
Taylor expansion above.  From this point on, \(N_m(0)\) denotes the
continuous extension supplied by Lemma~\ref{lem:Nm-removable-singularity}.
\end{remark}

\subsection{Whitened finite block}
\label{subsec:whitened-finite-block}

\begin{definition}[Whitened finite block]
\label{def:whitened-finite-block}
For \((m,s)\in\mathcal D_T\), with \(N_m(0)\) understood by continuous
extension, whenever \(G_{m,-}(s)\) and \(G_{m,+}(s)\) are positive
definite, let \(G_{m,\pm}(s)^{-1/2}\) denote the positive definite
inverse square root and set
\begin{equation}
\label{eq:omega-finite}
\widehat\Om_\z(s;m)
:= G_{m,-}(s)^{-1/2}\,N_m(s)\,G_{m,+}(s)^{-1/2}.
\end{equation}
\end{definition}

\begin{corollary}[Coalescence: \(\widehat\Om_\z(0;m) = I_2\)]
\label{cor:omega-coalescence}
For every \(m\in I_T\) at sufficiently large height,
\(\widehat\Om_\z(s;m)\) extends continuously to \(s=0\), along
\((m,s)\in\mathcal D_T\), with \(\widehat\Om_\z(0;m) = I_2\).
\end{corollary}

\begin{proof}
By Lemma~\ref{lem:same-point-gram-positivity}, \(J(m)\succ 0\) at
every sufficiently large retained midpoint \(m\).  By continuity,
\(G_{m,\pm}(s) = J(m\pm s/2)\) remain positive definite for all
sufficiently small \(|s|\) with \((m,s)\in\mathcal D_T\), and the map
\(G\mapsto G^{-1/2}\) is continuous on the positive definite cone.
Lemma~\ref{lem:Nm-removable-singularity} gives \(N_m(0) = J(m)\),
while \(G_{m,\pm}(0) = J(m)\) by
Lemma~\ref{lem:finite-blocks-exact-interpretation}.  Therefore
\(\widehat\Om_\z(0;m) = J(m)^{-1/2}\,J(m)\,J(m)^{-1/2} = I_2\).
\end{proof}

\subsection{Whitening loss}
\label{subsec:whitening-loss}

\begin{lemma}[Operator-norm bound for \(G_{m,\pm}^{-1/2}\)]
\label{lem:whitening-loss}
For \(\Ph_T = \theta\) on \(I_T\), after enlarging the global
lower-height cutoff if necessary, \(G_{m,\pm}(s)\) are positive
definite for every \((m,s)\in\mathcal D_T\), and there is a quantity
\(\varepsilon_T\to 0\), independent of \((m,s)\in\mathcal D_T\), such
that
\begin{equation}
\label{eq:whitening-loss-bound}
        \bigl\|G_{m,\pm}(s)^{-1/2}\bigr\|_{\op}
        \;\le\;
        \sqrt{\frac{\pi}{2\,q(t_\pm)}}\,(1+\varepsilon_T)
        \qquad (T\to\infty).
\end{equation}
After possibly enlarging the same lower-height cutoff,
\begin{equation}
\label{eq:whitening-loss-polynomial}
        \bigl\|G_{m,\pm}(s)^{-1/2}\bigr\|_{\op}
        \;\le\; Q(T)^{C_W}
\end{equation}
uniformly on \(\mathcal D_T\), with absolute exponent \(C_W = 0\)
whenever \(Q(T)\ge 1\).
\end{lemma}

\begin{proof}
By Lemma~\ref{lem:finite-blocks-exact-interpretation},
\(G_{m,\pm}(s) = J(t_\pm)\); since \((m,s)\in\mathcal D_T\) implies
\(t_\pm\in I_T\subset[T/2, 2T]\), the theta-derivative estimates of
Lemma~\ref{lem:theta-derivative-asymptotics} are uniform at
\(t = t_\pm\).  The proof of Lemma~\ref{lem:same-point-gram-positivity}
then applies with \(T\) replaced by an arbitrary \(t\in I_T\), giving,
uniformly on \(I_T\),
\[
        \lambda_{\min}(J(t))
        = \frac{2q(t)}{\pi}\,(1 + o(1)).
\]
Equivalently, for some \(\delta_T\to 0\) independent of \(t\in I_T\),
\[
        \lambda_{\min}(J(t))
        \;\ge\; \frac{2q(t)}{\pi}\,(1-\delta_T).
\]
For large \(T\), \(1-\delta_T > 0\); the finite same-point blocks are
positive definite, and
\[
        \bigl\|G_{m,\pm}(s)^{-1/2}\bigr\|_{\op}
        = \lambda_{\min}(G_{m,\pm}(s))^{-1/2}
        \;\le\; \sqrt{\frac{\pi}{2\,q(t_\pm)}}\,(1-\delta_T)^{-1/2}.
\]
Absorbing \((1-\delta_T)^{-1/2}\) into \(1+\varepsilon_T\) gives
\eqref{eq:whitening-loss-bound}.

For the polynomial bound, the uniform asymptotic above together with
Lemma~\ref{lem:phase-derivative-lower-bound} gives
\(\lambda_{\min}(J(t))\to\infty\) uniformly for \(t\in I_T\); after
enlarging the lower-height cutoff,
\(\lambda_{\min}(J(t))\ge 1\) uniformly on \(I_T\).  Hence
\(\|G_{m,\pm}(s)^{-1/2}\|_{\op}\le 1\le Q(T)^0\) whenever
\(Q(T)\ge 1\), which is \eqref{eq:whitening-loss-polynomial} with
\(C_W = 0\).
\end{proof}

\begin{remark}[Scope of Lemma~\ref{lem:whitening-loss}]
\label{rem:whitening-loss-scope}
Lemma~\ref{lem:whitening-loss} controls only the real-axis same-point
whitening factors attached to the theta phase.  It does not prove
actual-zeta transverse suppression, does not supply source-energy
input, does not give a uniform bound on \(\widehat\Om_\z(s;m)\) itself
beyond the whitening-factor loss, and does not give complex-disk or
holomorphic whitening estimates.  Any stronger whitening statement
needed for a later complex-analytic argument must be introduced as a
separate input or proved in a separate section.
\end{remark}
